\documentclass[../main.tex]{subfiles}
\begin{document}
本章主要参考\cite{hunter,brezis}。

% Reference: https://ocw.mit.edu/courses/18-303-linear-partial-differential-equations-fall-2006/pages/lecture-notes/
\section{基本类型与特征线}
\subsection{二阶方程分类与边界条件}\index{二阶方程分类与边界条件|hyperpage}
二维实系数方程的主部写为 $a u_{xx}+2b u_{xy}+c u_{yy}$，假定 $(a,b,c)\ne(0,0,0)$。逐点按
\[
b^2-ac<0\text{ 为椭圆型},\qquad b^2-ac=0\text{ 为抛物型},\qquad
b^2-ac>0\text{ 为双曲型}.
\]
典型方程为 $\Delta u=f$、$u_t=\kappa\Delta u$、$u_{tt}=c_0^2\Delta u$，其中 $\kappa,c_0>0$。
外单位法向为 $n$ 时，常见边界条件为
\[
\text{Dirichlet： }u=g,\qquad
\text{Neumann： }\partial_nu=g,\qquad
\text{Robin： }\alpha u+\beta\partial_nu=g.
\]
\subsection{一阶方程与特征线法}\index{一阶方程与特征线法|hyperpage}
对 $a(x,y,u)u_x+b(x,y,u)u_y=c(x,y,u)$，特征系统为
\[
\frac{\mathrm dx}{\mathrm ds}=a,\qquad
\frac{\mathrm dy}{\mathrm ds}=b,\qquad
\frac{\mathrm du}{\mathrm ds}=c.
\]
初值曲线 $(x_0(r),y_0(r),u_0(r))$ 在非特征条件 $a y_0'-b x_0'\ne0$ 下可局部确定解（假定系数与初值足够光滑）。
输运方程 $u_t+v u_x=0$ 的全直线初值解为 $u(x,t)=u_0(x-vt)$，其中 $v$ 为常数。
\section{热方程与波动方程}
\subsection{全空间热核与Duhamel公式}\index{全空间热核与Duhamel公式|hyperpage}
在 $\mathbb R^d$ 上，令 $G_t(x)=(4\pi\kappa t)^{-d/2}e^{-|x|^2/(4\kappa t)}$，则
\[
u_t-\kappa\Delta u=f,\quad u(\cdot,0)=u_0
\quad\Longrightarrow\quad
u(\cdot,t)=G_t*u_0+\int_0^tG_{t-s}*f(\cdot,s)\,\mathrm ds.
\]
该公式在数据光滑且有适当增长控制时给出经典解。对 $u_0\in L^p$、$1\le p<\infty$，齐次情形满足 $\|G_t*u_0-u_0\|_p\to0$；若 $u_0$ 有界一致连续，则一致恢复初值。
\subsection{一维波动方程的d'Alembert公式}\index{一维波动方程的d'Alembert公式|hyperpage}
对全直线初值问题 $u_{tt}=c^2u_{xx}$、$u(x,0)=\varphi(x)$、$u_t(x,0)=\psi(x)$，$c>0$，有
\[
u(x,t)=\frac{\varphi(x-ct)+\varphi(x+ct)}2
+\frac1{2c}\int_{x-ct}^{x+ct}\psi(s)\,\mathrm ds.
\]
例如 $\varphi\in C^2,\psi\in C^1$ 时给出经典解；该公式显示解只依赖区间 $[x-ct,x+ct]$ 内的初值。
\subsection{分离变量与齐次Dirichlet边界}\index{分离变量与齐次Dirichlet边界|hyperpage}
在 $0<x<L$ 上取 $u(0,t)=u(L,t)=0$，空间特征值与特征函数为
\[
\lambda_n=(n\pi/L)^2,\qquad X_n(x)=\sin(n\pi x/L),\quad n\ge1.
\]
热方程与波动方程的级数解分别为
\begin{align*}
u(x,t)&=\sum_{n=1}^\infty a_ne^{-\kappa\lambda_nt}\sin(n\pi x/L),\\
u(x,t)&=\sum_{n=1}^\infty\left[a_n\cos(\omega_nt)+\frac{b_n}{\omega_n}\sin(\omega_nt)\right]\sin(n\pi x/L),
\end{align*}
其中 $\omega_n=c\sqrt{\lambda_n}$，$a_n=(2/L)\int_0^Lu_0(x)\sin(n\pi x/L)\,\mathrm dx$，$b_n=(2/L)\int_0^Lv_0(x)\sin(n\pi x/L)\,\mathrm dx$。初值可按 $L^2$ 意义理解；逐项微分求经典解还需光滑性与边界相容条件。
\subsection{能量恒等式}\index{能量恒等式|hyperpage}
在有界光滑区域取齐次 Dirichlet 边界，且解足够光滑时，热方程满足
\[
\frac{\mathrm d}{\mathrm dt}\frac12\int_\Omega u^2\,\mathrm dx
=-\kappa\int_\Omega|\nabla u|^2\,\mathrm dx.
\]
齐次波动方程的能量守恒：
\[
E(t)=\frac12\int_\Omega(u_t^2+c^2|\nabla u|^2)\,\mathrm dx,\qquad E'(t)=0.
\]
边界非齐次或存在源项时，应保留边界通量或外力做功项。
\section{Laplace与Poisson方程}
\subsection{调和函数、最大值原理与Poisson核}\index{调和函数、最大值原理与Poisson核|hyperpage}
若 $\Delta u=0$，称 $u$ 为调和函数。球 $\overline{B_R(x_0)}$ 位于调和区域内部时，
\[
u(x_0)=\frac1{|\partial B_R|}\int_{\partial B_R(x_0)}u\,\mathrm dS
=\frac1{|B_R|}\int_{B_R(x_0)}u\,\mathrm dx.
\]
在有界区域上，$u\in C^2(\Omega)\cap C(\overline\Omega)$ 调和时最大最小值在边界取得；若连通且内部取最大值，则为常数。
单位圆盘上的连续边界数据 $g$ 给出 Poisson 积分
\[
u(r,\theta)=\frac1{2\pi}\int_0^{2\pi}
\frac{1-r^2}{1-2r\cos(\theta-\phi)+r^2}\,g(\phi)\,\mathrm d\phi,
\quad0\le r<1.
\]
\subsection{基本解与Neumann相容性}\index{基本解与Neumann相容性|hyperpage}
采用 $-\Delta\Phi=\delta_0$ 的符号约定，$d\ge2$ 的全空间基本解为
\[
\Phi(x)=\begin{cases}
-\dfrac1{2\pi}\ln|x|,&d=2,\\
\dfrac{|x|^{2-d}}{(d-2)|S^{d-1}|},&d\ge3.
\end{cases}
\]
等式按分布意义成立；对光滑紧支撑 $f$，$u=\Phi*f$ 满足 $-\Delta u=f$。
对有界连通光滑区域的 Neumann 问题 $-\Delta u=f$、$\partial_nu=g$，必要相容条件为
\[
\int_\Omega f\,\mathrm dx+\int_{\partial\Omega}g\,\mathrm dS=0.
\]
解至多相差常数，常用 $\int_\Omega u=0$ 固定这一自由度。
\section{弱导数与变分形式}
\subsection{Sobolev空间与弱解}\index{Sobolev空间与弱解|hyperpage}
局部可积 $u$ 的弱导数 $D^\alpha u=v$ 指对任意 $\varphi\in C_c^\infty(\Omega)$ 有
\[
\int_\Omega uD^\alpha\varphi\,\mathrm dx=(-1)^{|\alpha|}\int_\Omega v\varphi\,\mathrm dx.
\]
$W^{k,p}(\Omega)$ 由 $k$ 阶以内弱导数均在 $L^p$ 中的函数组成；$H^1=W^{1,2}$，
\[
\|u\|_{H^1}^2=\int_\Omega(|u|^2+|\nabla u|^2)\,\mathrm dx,\qquad
H_0^1=\overline{C_c^\infty(\Omega)}^{\,H^1}.
\]
在有界 Lipschitz 区域上，$u\in H_0^1$ 满足 Poincaré 不等式 $\|u\|_2\le C_\Omega\|\nabla u\|_2$。对 $f\in L^2$，Dirichlet 问题 $-\Delta u=f$ 的弱形式为
\[
\int_\Omega\nabla u\cdot\nabla v\,\mathrm dx=\int_\Omega fv\,\mathrm dx
\quad(\forall v\in H_0^1),\qquad u\in H_0^1.
\]
由 Lax--Milgram 定理有唯一弱解；它也唯一极小化 $J(v)=\tfrac12\int|\nabla v|^2-\int fv$。
% Further references: https://web.stanford.edu/class/math220a/ ; https://www.math.ucdavis.edu/~hunter/pdes/
\section{椭圆方程与边界积分补充}
\subsection{Green恒等式与Green函数表示}\index{Green恒等式与Green函数表示|hyperpage}
对光滑有界域及光滑函数，外法向约定给出
\begin{align*}
\int_\Omega(\nabla u\cdot\nabla v+u\Delta v)&=\int_{\partial\Omega}u\partial_nv,\\
\int_\Omega(u\Delta v-v\Delta u)&=\int_{\partial\Omega}(u\partial_nv-v\partial_nu).
\end{align*}
Dirichlet Green 函数满足 $-\Delta_yG(x,y)=\delta_x(y)$、边界上 $G=0$。对 $-\Delta u=f$、$u|_{\partial\Omega}=g$，
\[
u(x)=\int_\Omega G(x,y)f(y)\,\mathrm dy
-\int_{\partial\Omega}g(y)\partial_{n_y}G(x,y)\,\mathrm dS_y.
\]
负号来自本章采用 $-\Delta$ 的基本解约定；对称域问题有 $G(x,y)=G(y,x)$。
\subsection{球的Poisson核与Harnack估计}\index{球的Poisson核与Harnack估计|hyperpage}
球 $B_R(0)\subset\mathbb R^d$ 上调和 Dirichlet 解为
\[
u(x)=\int_{|y|=R}\frac{R^2-|x|^2}{|S^{d-1}|R|x-y|^d}g(y)\,\mathrm dS_y.
\]
若 $u\ge0$ 在 $B_{2r}$ 调和，则 $\sup_{B_r}u\le C_d\inf_{B_r}u$。调和函数的内点导数满足
\[
|D^\alpha u(x)|\le C_{d,\alpha}r^{-|\alpha|}\sup_{B_r(x)}|u|.
\]
若 $\Delta u\ge0$，弱最大值原理给出 $\sup_\Omega u\le\sup_{\partial\Omega}u$；应用于一般椭圆算子时需检查主部正定性和零阶项符号。
\subsection{一致椭圆性、弱能量与正则性}\index{一致椭圆性、弱能量与正则性|hyperpage}
考虑 $Lu=-\operatorname{div}(A\nabla u)+cu$，其中 $A$ 有界且 $\xi^{\mathsf T}A(x)\xi\ge\lambda|\xi|^2$、$c\ge0$。齐次 Dirichlet 解满足
\[
\lambda\|\nabla u\|_2^2\le\langle f,u\rangle,\qquad
\|u\|_{H_0^1}\le C\|f\|_{H^{-1}}.
\]
对 $Lu=0$ 的弱解和紧支撑截断 $\eta$，Caccioppoli 估计为 $\int\eta^2|\nabla u|^2\le C\int|\nabla\eta|^2|u|^2$。若 $\Omega'\Subset\Omega$、$-\Delta u=f$，则
\begin{align*}
\|u\|_{W^{2,p}(\Omega')}&\le C(\|f\|_{L^p(\Omega)}+\|u\|_{L^p(\Omega)})\quad(1<p<\infty),\\
\|u\|_{C^{2,\alpha}(\Omega')}&\le C(\|f\|_{C^{0,\alpha}(\Omega)}+\|u\|_{C^0(\Omega)})\quad(0<\alpha<1).
\end{align*}
这是内点估计；全局估计还依赖边界光滑性、边值及相容条件，粗糙系数不能直接套用。
\subsection{特征函数展开与Rayleigh商}\index{特征函数展开与Rayleigh商|hyperpage}
有界光滑域的 Dirichlet Laplace 算子有 $L^2$ 正交特征基 $-\Delta\phi_n=\lambda_n\phi_n$，$0<\lambda_1\le\lambda_2\to\infty$。有
\[
\lambda_k=\min_{\substack{V\subset H_0^1\\\dim V=k}}\max_{0\ne u\in V}
\frac{\int|\nabla u|^2}{\int|u|^2},\qquad
u=\sum_n\frac{\langle f,\phi_n\rangle}{\lambda_n}\phi_n\quad(-\Delta u=f).
\]
热方程对应 $e^{-\kappa\lambda_nt}$，波动方程对应 $\cos(c\sqrt{\lambda_n}t)$ 和 $\sin(c\sqrt{\lambda_n}t)/(c\sqrt{\lambda_n})$。Weyl 渐近公式为
\[
\#\{n:\lambda_n\le\Lambda\}\sim\frac{\omega_d|\Omega|}{(2\pi)^d}\Lambda^{d/2},
\]
其中 $\omega_d$ 为单位球体积，按重数计数。
\section{演化方程的核、估计与守恒}
\subsection{热方程平滑化与比较原理}\index{热方程平滑化与比较原理|hyperpage}
全空间热半群对 $1\le p\le q\le\infty$ 满足
\[
\|D^\alpha(G_t*f)\|_q
\le C(\kappa t)^{-|\alpha|/2-d(1/p-1/q)/2}\|f\|_p.
\]
非负初值保持非负，$L^p$ 范数不增，且 $\int G_t*f=\int f$（$f\in L^1$）。有界域上 $u_t-\kappa\Delta u\le0$ 的光滑解最大值由抛物边界 $\partial\Omega\times[0,T]\cup\overline\Omega\times\{0\}$ 控制。齐次 Dirichlet 边界给出
\[
\|u(t)\|_2\le e^{-\kappa\lambda_1t}\|u_0\|_2.
\]
\subsection{多维波动公式与有限传播}\index{多维波动公式与有限传播|hyperpage}
三维全空间初值 $u(0)=\varphi,u_t(0)=\psi$ 的 Kirchhoff 公式（$t>0$）为
\[
u(x,t)=\partial_t\left[\frac1{4\pi c^2t}\int_{|y-x|=ct}\varphi(y)\,\mathrm dS_y\right]
+\frac1{4\pi c^2t}\int_{|y-x|=ct}\psi(y)\,\mathrm dS_y.
\]
对任意维数，Fourier 形式为
\[
\widehat u(\xi,t)=\cos(ct|\xi|)\widehat\varphi(\xi)
+\frac{\sin(ct|\xi|)}{c|\xi|}\widehat\psi(\xi),
\]
在 $\xi=0$ 用连续值 $t$。紧支撑初值的支撑以速度至多 $c$ 扩张；奇数维三维的球面依赖不应推广成所有维数的球面公式。
\subsection{Schrödinger方程与色散}\index{Schrödinger方程与色散|hyperpage}
采用 $iu_t+\Delta u=0$，有 $\widehat u(t,\xi)=e^{-it|\xi|^2}\widehat u_0(\xi)$，且对 $t\ne0$，
\[
u(x,t)=(4\pi it)^{-d/2}\int e^{i|x-y|^2/(4t)}u_0(y)\,\mathrm dy,
\qquad\|u(t)\|_\infty\le(4\pi|t|)^{-d/2}\|u_0\|_1.
\]
相位因子按连续 Fourier 核约定解释。$\|u(t)\|_2=\|u_0\|_2$，光滑数据还守恒 $\|\nabla u(t)\|_2$。有源问题 $iu_t+\Delta u=F$ 的 Duhamel 项为 $-i\int_0^te^{i(t-s)\Delta}F(s)\,\mathrm ds$。
\section{非线性方程与弱解补充}
\subsection{守恒律、间断速度与熵条件}\index{守恒律、间断速度与熵条件|hyperpage}
一维标量守恒律 $u_t+f(u)_x=0$ 的光滑特征满足 $x'=f'(u)$、$u'=0$。间断两侧状态 $u_L,u_R$ 的 Rankine--Hugoniot 速度为
\[
s=\frac{f(u_R)-f(u_L)}{u_R-u_L}.
\]
对凸通量，压缩激波满足 $f'(u_L)>s>f'(u_R)$。Kružkov 熵不等式为（任意常数 $k$）
\[
\partial_t|u-k|+\partial_x[\operatorname{sgn}(u-k)(f(u)-f(k))]\le0
\]
按分布成立。对 $f(u)=u^2/2$，$u_L>u_R$ 的激波速度为 $(u_L+u_R)/2$；若 $u_L<u_R$，稀疏波为 $u=u_L$（$x/t\le u_L$）、$u=x/t$（$u_L<x/t<u_R$）、$u=u_R$（$x/t\ge u_R$）。
\subsection{Burgers黏性正则化与Cole--Hopf}\index{Burgers黏性正则化与Cole-Hopf|hyperpage}
\[
u_t+uu_x=\nu u_{xx},\qquad u=-2\nu\partial_x\log\theta,\qquad
\theta_t=\nu\theta_{xx}\quad(\nu>0,\ \theta>0).
\]
由初值 $u_0$ 可取 $\theta_0(x)=\exp[-(2\nu)^{-1}\int_0^xu_0(s)\,\mathrm ds]$，再用热核求 $\theta$；需保证相关积分有意义。无黏光滑解在特征交叉前为 $u(x,t)=u_0(\xi)$、$x=\xi+tu_0(\xi)$，梯度为
\[
u_x(x,t)=\frac{u_0'(\xi)}{1+tu_0'(\xi)}.
\]
若 $\inf u_0'<0$ 且适当有界，最早梯度爆破时间为 $-1/\inf u_0'$。
\subsection{Hamilton--Jacobi与Hopf--Lax公式}\index{Hamilton-Jacobi与Hopf-Lax公式|hyperpage}
对凸超线性 Hamiltonian $H$，令 $L(v)=\sup_p\{p\cdot v-H(p)\}$。在适当 Lipschitz 初值下，$u_t+H(\nabla u)=0$ 的黏性解为
\[
u(x,t)=\inf_y\left\{u_0(y)+tL\left(\frac{x-y}t\right)\right\}\quad(t>0).
\]
例如 $H(p)=|p|^2/2$ 时 $L(v)=|v|^2/2$。黏性次解条件是在测试函数 $\phi$ 与 $u$ 相切于局部极大处要求 $\phi_t+H(\nabla\phi)\le0$；超解在局部极小处反向，不能把黏性解理解为处处经典可微。
\subsection{p-Laplace与变分结构}\index{p-Laplace与变分结构|hyperpage}
对 $1<p<\infty$，$\Delta_pu=\operatorname{div}(|\nabla u|^{p-2}\nabla u)$。能量
\[
J(u)=\frac1p\int_\Omega|\nabla u|^p-\langle f,u\rangle
\]
在 $W_0^{1,p}$ 上的极小点满足 $\int|\nabla u|^{p-2}\nabla u\cdot\nabla v=\langle f,v\rangle$。严格凸性及 Poincaré 不等式给出存在唯一性（$f\in W^{-1,p'}$）。更一般的一阶变分为
\[
J(u)=\int F(x,u,\nabla u)\quad\Longrightarrow\quad
F_u-\operatorname{div}F_{\nabla u}=0
\]
在相应弱或经典意义下成立。
\subsection{不可压流体的基本恒等式}\index{不可压流体的基本恒等式|hyperpage}
Navier--Stokes 方程写为
\[
u_t+(u\cdot\nabla)u+\nabla p=\nu\Delta u+f,\qquad\nabla\cdot u=0.
\]
在周期域、全空间充分衰减或无滑移边界下，光滑解满足
\[
\frac12\frac{\mathrm d}{\mathrm dt}\|u\|_2^2+\nu\|\nabla u\|_2^2=\int f\cdot u.
\]
无外力三维涡量 $\omega=\nabla\times u$ 满足 $\omega_t+(u\cdot\nabla)\omega=(\omega\cdot\nabla)u+\nu\Delta\omega$；二维标量涡量没有拉伸项。能量估计本身不保证三维任意光滑初值的全局光滑解。

% 第二次扩充（第1组）：椭圆理论。
% References: https://www.math.cmu.edu/~gautam/c/2026-632/notes/elliptic-equations.html
% https://www.math.ucdavis.edu/~hunter/pdes/pde_notes.pdf
\section{椭圆方程的边界与粗糙系数理论}
本轮补充中的常数 $C$ 随维数、椭圆性、指数及指定区域而变，通常不随被估计的解变化。Sobolev 嵌入、奇异积分与变分工具分别见\hyperref[chap:real-analysis]{实分析}、\hyperref[chap:harmonic-analysis]{调和分析}和\hyperref[chap:calculus-of-variations]{变分法}。
\subsection{Hopf边界引理与法向符号}\index{Hopf边界引理与法向符号|hyperpage}
设 $\Omega$ 连通，边界在 $x_0\in\partial\Omega$ 为 $C^1$ 且满足内切球条件。若 $u\in C^2(\Omega)\cap C^1(\overline\Omega)$、$\Delta u\ge0$，且非常数 $u$ 在 $x_0$ 达到全域最大值，则对外法向有
\[
\partial_nu(x_0)>0.
\]
若 $\Delta u\le0$ 且在边界取非常数最小值，则 $\partial_nu(x_0)<0$。证明常用环域屏障
\[
w(x)=e^{-a|x-z|^2}-e^{-aR^2},\qquad
\Delta w=(4a^2|x-z|^2-2ad)e^{-a|x-z|^2},
\]
其中 $B_R(z)$ 为内切球，选足够大的 $a$。一般一致椭圆算子还须检查低阶项；边界尖点处不能直接套用内切球结论。
\subsection{ABP最大值估计与接触集}\index{ABP最大值估计与接触集|hyperpage}
有界域 $\Omega\subset\mathbb R^d$ 上，令 $Lu=a^{ij}(x)D_{ij}u$，$A$ 对称且 $\lambda I\le A\le\Lambda I$。对 $u\in C^2(\Omega)\cap C(\overline\Omega)$、$Lu\ge f$，ABP 估计给出
\[
\sup_\Omega u\le\sup_{\partial\Omega}u^+
+C(d,\lambda)\operatorname{diam}(\Omega)\|f^-\|_{L^d(\Omega)}.
\]
其中 $u^+=\max(u,0)$、$f^-=\max(-f,0)$。更精细的版本只积分上接触集：在这些点，图像有从上方支撑它的仿射超平面。若改用 $-a^{ij}D_{ij}$，右端正负部也须相应改变。带漂移、零阶项及弱正则性的版本需要另外列出其假设。
\subsection{De Giorgi--Nash--Moser估计}\index{De Giorgi-Nash-Moser估计|hyperpage}
设 $u$ 为 $-\operatorname{div}(A(x)\nabla u)=0$ 在 $B_{2R}$ 上的弱解，$A$ 仅需有界可测，满足 $\xi^{\mathsf T}A\xi\ge\lambda|\xi|^2$ 及 $|A\xi|\le\Lambda|\xi|$。存在 $\alpha\in(0,1)$，使
\[
[u]_{C^{0,\alpha}(B_R)}
\le CR^{-\alpha}\left(R^{-d}\int_{B_{2R}}|u|^2\,\mathrm dx\right)^{1/2}.
\]
若 $u\ge0$，还有 $\operatorname*{ess\,sup}_{B_R}u\le C\operatorname*{ess\,inf}_{B_R}u$。这里的结论是解本身 Hölder 连续；仅凭可测系数不能推出梯度连续或二阶导数的 Schauder 估计。该标量理论不能无条件推广到一般椭圆系统。
\subsection{全局边界正则性与椭圆估计}\index{全局边界正则性与椭圆估计|hyperpage}
有界 $C^{1,1}$ 域上，齐次 Dirichlet 问题 $-\Delta u=f$ 满足
\[
\|u\|_{W^{2,p}(\Omega)}\le C\|f\|_{L^p(\Omega)},\qquad1<p<\infty.
\]
这包括 $u\in W^{2,p}(\Omega)\cap W_0^{1,p}(\Omega)$ 的存在唯一性。若域为 $C^{2,\alpha}$、$f\in C^{0,\alpha}(\overline\Omega)$，且边值为某 $g\in C^{2,\alpha}(\overline\Omega)$ 的迹，则
\[
\|u\|_{C^{2,\alpha}(\overline\Omega)}
\le C\bigl(\|f\|_{C^{0,\alpha}(\overline\Omega)}
+\|g\|_{C^{2,\alpha}(\overline\Omega)}\bigr).
\]
角域、混合边界及粗糙系数需重新分析奇性；不能把内点估计直接当作全局边界估计。

% 第二次扩充（第2组）：抛物理论。
% References: https://virtualmath1.stanford.edu/~ryzhik/STANFORD/STANF272-21/lecture-notes-stanf272-21.pdf
% https://arxiv.org/abs/2003.04554
\section{抛物方程的弱解、正则性与行波}
\subsection{弱抛物解与能量空间}\index{弱抛物解与能量空间|hyperpage}
有界 Lipschitz 域上取 $V=H_0^1(\Omega)$、$H=L^2(\Omega)$，$V^*=H^{-1}(\Omega)$。若 $u\in L^2(0,T;V)$、$u_t\in L^2(0,T;V^*)$，则 $u$ 有 $C([0,T];H)$ 代表，且
\[
\frac{\mathrm d}{\mathrm dt}\frac12\|u(t)\|_H^2
=\langle u_t(t),u(t)\rangle_{V^*,V}\quad\text{a.e.}
\]
对 $u_t-\operatorname{div}(A\nabla u)=f$、齐次 Dirichlet 边界，若 $A$ 有界、可测且一致椭圆，$u_0\in H$、$f\in L^2(0,T;V^*)$，则弱解满足
\[
\sup_{0\le t\le T}\|u(t)\|_2^2+\int_0^T\|\nabla u\|_2^2\,\mathrm dt
\le C\left(\|u_0\|_2^2+\int_0^T\|f\|_{V^*}^2\,\mathrm dt\right).
\]
时间导数在对偶空间中理解；不能要求任意能量弱解都具有逐点经典时间导数。
\subsection{抛物型Harnack不等式与时间方向}\index{抛物型Harnack不等式与时间方向|hyperpage}
设 $u\ge0$ 弱解 $u_t-\operatorname{div}(A(x,t)\nabla u)=0$，定义
\begin{gather*}
Q=B_{2r}(x_0)\times(t_0-4r^2,t_0),\\
Q_-=B_r(x_0)\times(t_0-3r^2,t_0-2r^2),\qquad
Q_+=B_r(x_0)\times(t_0-r^2,t_0).
\end{gather*}
在 $Q$ 内系数有界可测且一致椭圆时，
\[
\operatorname*{ess\,sup}_{Q_-}u\le C\operatorname*{ess\,inf}_{Q_+}u.
\]
较早时刻的上界由较晚时刻的下界控制，不能交换两个时间柱。相应局部正则性采用抛物距离 $|x-y|+|t-s|^{1/2}$，即空间 Hölder 指数 $\alpha$ 对应时间指数 $\alpha/2$。
\subsection{热方程最大正则性}\index{热方程最大正则性|hyperpage}
在 $\mathbb R^d\times(0,T)$ 上，取 $u_t-\kappa\Delta u=f$、$u(0)=0$、$\kappa>0$。对 $1<p,q<\infty$，热核解满足
\[
\|u_t\|_{L^q(0,T;L^p)}
+\kappa\|D^2u\|_{L^q(0,T;L^p)}
\le C_{d,p,q}\|f\|_{L^q(0,T;L^p)}.
\]
非零初值需要相应的迹空间，不能在右端直接删去初值项。一个便于使用的 Hilbert 情形是在有界 $C^{1,1}$ 域取齐次 Dirichlet 边界、$u_0\in H_0^1$、$f\in L^2(0,T;L^2)$，此时
\[
\|u_t\|_{L^2(0,T;L^2)}+\|u\|_{L^2(0,T;H^2)}
\le C\bigl(\|f\|_{L^2(0,T;L^2)}+\|u_0\|_{H^1}\bigr).
\]
端点 $p=1,\infty$ 或 $q=1,\infty$ 一般不能沿用上述强范数结论。
\subsection{反应扩散、比较原理与KPP行波}\index{反应扩散、比较原理与KPP行波|hyperpage}
对 $u_t=D\Delta u+F(u)$、$D>0$，局部 Lipschitz 的 $F$ 允许在适当有界解类中用比较原理控制解；若 $F(0)=F(1)=0$，则初边值在 $[0,1]$ 内时解保持在该区间。一维 Fisher--KPP 方程
\[
u_t=Du_{xx}+r u(1-u),\qquad D,r>0,
\]
取 $u(x,t)=U(x-ct)$，得到
\[
DU''+cU'+rU(1-U)=0,\quad U(-\infty)=1,\quad U(+\infty)=0.
\]
从 $1$ 向右传播至 $0$ 的单调行波存在当且仅当 $c\ge c_*=2\sqrt{Dr}$。前缘 $U\sim e^{-\lambda z}$ 的线性化满足 $D\lambda^2-c\lambda+r=0$；其他反应项不一定具有同样的最小速度公式。

% 第二次扩充（第3组）：波与双曲系统。
% References: https://www.math.ucdavis.edu/~hunter/pdes/pde_notes.pdf
% https://math.ucla.edu/~tao/preprints/chapter.pdf
\section{多维波、局部能量与双曲系统}
\subsection{二维波动方程的Poisson公式}\index{二维波动方程的Poisson公式|hyperpage}
对 $\mathbb R^2$ 上 $u_{tt}-c^2\Delta u=0$、$u(0)=\varphi$、$u_t(0)=\psi$，$c>0$，定义
\[
\mathcal P_t h(x)=\frac1{2\pi c}\int_{|y-x|<ct}
\frac{h(y)}{\sqrt{c^2t^2-|y-x|^2}}\,\mathrm dy.
\]
则对足够光滑初值和 $t>0$，
\[
u(x,t)=\partial_t\mathcal P_t\varphi(x)+\mathcal P_t\psi(x).
\]
积分依赖整个圆盘，而非仅圆周。此式可由三维 Kirchhoff 公式降维得到；换元 $y=x+ctz$ 后可在固定单位圆盘中处理时间微分。核在圆周处的奇性可积。
\subsection{局部能量通量与依赖区域}\index{局部能量通量与依赖区域|hyperpage}
对带源波动方程 $u_{tt}-c^2\Delta u=f$，令 $e=(u_t^2+c^2|\nabla u|^2)/2$，有局部守恒式
\[
\partial_te+\operatorname{div}(-c^2u_t\nabla u)=fu_t.
\]
在固定光滑子域 $D$ 上，
\[
\frac{\mathrm d}{\mathrm dt}\int_De
=c^2\int_{\partial D}u_t\partial_nu\,\mathrm dS+\int_Dfu_t\,\mathrm dx.
\]
对收缩球 $D_t=B_{R-ct}(x_0)$、$0<t<R/c$，运动边界项非正，因此
\[
\frac{\mathrm d}{\mathrm dt}\int_{D_t}e\le\int_{D_t}fu_t.
\]
这给出有限传播与局部唯一性；零源且初值在 $B_R$ 内均为零时，解在相应前向收缩锥内为零。不能把扩张球误当作同一单调能量区域。
\subsection{对称双曲系统与Sobolev能量}\index{对称双曲系统与Sobolev能量|hyperpage}
考虑 $\partial_tU+\sum_{j=1}^dA_j(x,t)\partial_jU+B(x,t)U=F$，其中实矩阵 $A_j=A_j^{\mathsf T}$。在全空间衰减或周期条件下，光滑解满足
\begin{align*}
\frac12\frac{\mathrm d}{\mathrm dt}\|U\|_2^2
&=\int U\cdot F\,\mathrm dx\\
&\quad+\frac12\int U^{\mathsf T}
\left(\sum_j\partial_jA_j-B-B^{\mathsf T}\right)U\,\mathrm dx.
\end{align*}
对整数 $s\ge0$，系数及其所需空间导数有界时，交换微分并用 Gronwall 得
\[
\|U(t)\|_{H^s}\le e^{Ct}
\left(\|U(0)\|_{H^s}+\int_0^t\|F(\tau)\|_{H^s}\,\mathrm d\tau\right).
\]
取 $C$ 控制 $B$ 的 $W^{s,\infty}$ 范数及 $A_j$ 的 $W^{s+1,\infty}$ 范数即可。有边界时须另外控制入射特征与边界能量通量。
\subsection{Klein--Gordon传播子与能量}\index{Klein-Gordon传播子与能量|hyperpage}
全空间方程 $u_{tt}-c^2\Delta u+m^2u=F$，其中 $m,c>0$。记 $\omega(\xi)=\sqrt{c^2|\xi|^2+m^2}$，则
\begin{align*}
\widehat u(t,\xi)&=\cos(t\omega)\widehat\varphi(\xi)
+\frac{\sin(t\omega)}{\omega}\widehat\psi(\xi)\\
&\quad+\int_0^t\frac{\sin((t-s)\omega)}{\omega}\widehat F(s,\xi)\,\mathrm ds.
\end{align*}
齐次情形守恒正能量
\[
E_{\rm KG}(t)=\frac12\int\bigl(|u_t|^2+c^2|\nabla u|^2+m^2|u|^2\bigr)\,\mathrm dx.
\]
质量项给出额外 $L^2$ 控制，有限传播速度仍为 $c$；全空间色散估计须按频率及正则性另行说明，不能只由能量守恒推出。

% 第二次扩充（第4组）：色散与NLS。
% Reference: https://math.ucla.edu/~tao/preprints/chapter.pdf
\section{色散时空估计与非线性Schrödinger方程}
\subsection{Strichartz估计与允许指数}\index{Strichartz估计与允许指数|hyperpage}
在 $\mathbb R^d$ 上，称 $(q,r)$ 为 Schrödinger 允许对，若
\[
2\le q,r\le\infty,\qquad\frac2q+\frac dr=\frac d2,\qquad
(q,r,d)\ne(2,\infty,2).
\]
对允许对 $(q,r)$、$(\widetilde q,\widetilde r)$，有
\begin{gather*}
\|e^{it\Delta}f\|_{L_t^qL_x^r(\mathbb R\times\mathbb R^d)}\le C\|f\|_2,\\
\left\|\int_0^te^{i(t-s)\Delta}F(s)\,\mathrm ds\right\|_{L^q(0,T;L^r)}
\le C\|F\|_{L^{\widetilde q'}(0,T;L^{\widetilde r'})}.
\end{gather*}
撇号表示 Hölder 共轭。这里是欧氏全空间、常系数自由传播子的估计；有界域、势场或一般流形不一定保持同一全局结论，二维被排除的端点也不能补回。
\subsection{非线性Schrödinger的积分形式与守恒量}\index{非线性Schrödinger的积分形式与守恒量|hyperpage}
固定符号约定 $iu_t+\Delta u=\mu|u|^{p-1}u$，$p>1$、$\mu\in\{-1,1\}$，则
\[
u(t)=e^{it\Delta}u_0-i\mu\int_0^te^{i(t-s)\Delta}
\bigl(|u(s)|^{p-1}u(s)\bigr)\,\mathrm ds.
\]
对足够正则、衰减且能量有限的解，质量及能量守恒：
\[
M(u)=\int|u|^2,\qquad
E(u)=\frac12\int|\nabla u|^2+\frac\mu{p+1}\int|u|^{p+1}.
\]
$\mu=1$ 为散焦、$\mu=-1$ 为聚焦。积分方程用于固定点构造；$H^1$ 解的能量定义还要求 $u\in L^{p+1}$，不能对任意指数默认该项有限。
\subsection{尺度临界性与Gagliardo--Nirenberg控制}\index{尺度临界性与Gagliardo-Nirenberg控制|hyperpage}
上述 NLS 的尺度变换及临界 Sobolev 指数为
\[
u_\lambda(x,t)=\lambda^{2/(p-1)}u(\lambda x,\lambda^2t),\qquad
s_c=\frac d2-\frac2{p-1},
\]
\[
\|u_\lambda(0)\|_{\dot H^s}=\lambda^{s-s_c}\|u_0\|_{\dot H^s}.
\]
质量临界为 $p=1+4/d$；$d\ge3$ 时能量临界为 $p=1+4/(d-2)$。对 $u\in H^1(\mathbb R^d)$，
\[
\|u\|_{p+1}\le C\|\nabla u\|_2^\theta\|u\|_2^{1-\theta},\qquad
\theta=d\left(\frac12-\frac1{p+1}\right).
\]
这里 $d\ge3$ 时 $2\le p+1\le2d/(d-2)$，$d=1,2$ 时可取任意有限 $p+1\ge2$。临界性描述尺度，不单独给出适定性或散射结论。
\subsection{Virial恒等式与聚焦爆破判据}\index{Virial恒等式与聚焦爆破判据|hyperpage}
假设解足够光滑、衰减，且 $xu\in L^2$。方差 $V(t)=\int|x|^2|u|^2$ 满足
\begin{align*}
V'(t)&=4\operatorname{Im}\int x\cdot\overline u\nabla u,\\
V''(t)&=8\int|\nabla u|^2
+\frac{4\mu d(p-1)}{p+1}\int|u|^{p+1}\\
&=16E(u)+\frac{4\mu[d(p-1)-4]}{p+1}\int|u|^{p+1}.
\end{align*}
质量临界时 $V''=16E$。聚焦且 $p\ge1+4/d$、$E(u_0)<0$ 时，有 $V''\le16E(u_0)<0$；在满足上述恒等式与延拓理论的解类中，这排除正时间全局正则解。有限方差、正则性及非线性指数条件不可删除；局部化 virial 需要保留截断误差项。

% 第二次扩充（第5组）：一维标量守恒律。
% References: https://web.stanford.edu/class/math220a/handouts/conservation.pdf
% https://andreask.cs.illinois.edu/dl/academic/notes/257/257.pdf
\section{标量守恒律的熵、收缩与黏性极限}
\subsection{一般熵对与积分熵条件}\index{一般熵对与积分熵条件|hyperpage}
对一维标量方程 $u_t+f(u)_x=0$，熵--熵通量对 $(\eta,q)$ 满足
\[
\eta''\ge0,\qquad q'(u)=\eta'(u)f'(u).
\]
光滑解满足 $\partial_t\eta(u)+\partial_xq(u)=0$；熵弱解要求该量在分布意义下 $\le0$。对任意非负 $\phi\in C_c^\infty(\mathbb R\times[0,\infty))$，写为
\[
\int_0^\infty\!\int_{\mathbb R}
\bigl(\eta(u)\phi_t+q(u)\phi_x\bigr)\,\mathrm dx\,\mathrm dt
+\int_{\mathbb R}\eta(u_0)\phi(x,0)\,\mathrm dx\ge0.
\]
Kružkov 族为 $\eta_k(u)=|u-k|$、$q_k(u)=\operatorname{sgn}(u-k)(f(u)-f(k))$。时间积分分部带来的初值项与不等号方向必须一致。
\subsection{L1收缩、比较与局部依赖}\index{L1收缩、比较与局部依赖|hyperpage}
若 $f\in C^1$，$u,v$ 是全直线上有界熵解，初值差在 $L^1$ 中，则
\[
\|u(t)-v(t)\|_{L^1(\mathbb R)}\le\|u_0-v_0\|_{L^1(\mathbb R)}.
\]
这给出熵解的唯一性与对初值的稳定性；还满足 $u_0\le v_0\Rightarrow u(t)\le v(t)$。若两解的值位于区间 $I$，取 $M=\sup_{z\in I}|f'(z)|<\infty$，则
\[
\int_{|x-x_0|<R}|u(x,t)-v(x,t)|\,\mathrm dx
\le\int_{|x-x_0|<R+Mt}|u_0(x)-v_0(x)|\,\mathrm dx.
\]
因此差异的传播速度由 $M$ 控制；仅满足分布方程而不满足熵条件的弱解一般不具备该唯一性结论。
\subsection{Oleinik单边估计与特征交叉}\index{Oleinik单边估计与特征交叉|hyperpage}
若 $f\in C^2$ 且在解的有界值域上 $f''\ge a>0$，熵解满足对每个 $t>0$、几乎处处 $x<y$，
\[
u(y,t)-u(x,t)\le\frac{y-x}{at}.
\]
这是单边 Lipschitz 控制，允许向下跳跃而限制过陡的向上增长。对光滑初值，在特征交叉前有
\[
x=\xi+t f'(u_0(\xi)),\qquad
u_x(x,t)=\frac{u_0'(\xi)}{1+t f''(u_0(\xi))u_0'(\xi)}.
\]
当 $m_0=\inf_\xi[f''(u_0(\xi))u_0'(\xi)]<0$，并具备足够有界性保证特征构造时，首次梯度奇性的时间为 $T_*=-1/m_0$。该时间属于光滑解，不是熵解存在性的终点。
\subsection{消失黏性、全变差与熵耗散}\index{消失黏性、全变差与熵耗散|hyperpage}
对 $u_t^\varepsilon+f(u^\varepsilon)_x=\varepsilon u_{xx}^\varepsilon$，光滑凸熵满足
\[
\partial_t\eta(u^\varepsilon)+\partial_xq(u^\varepsilon)
=\varepsilon\partial_{xx}\eta(u^\varepsilon)
-\varepsilon\eta''(u^\varepsilon)|u_x^\varepsilon|^2.
\]
若 $u_0\in L^1\cap L^\infty\cap BV(\mathbb R)$，用保持统一界的光滑逼近取初值，$f\in C^2$，则有
\[
\|u^\varepsilon(t)\|_\infty\le\|u_0^\varepsilon\|_\infty,\qquad
\operatorname{TV}(u^\varepsilon(t))\le\operatorname{TV}(u_0^\varepsilon).
\]
当初值在 $L^1$ 收敛时，$u^\varepsilon$ 收敛到唯一熵解，可在每个有限时间段按 $C([0,T];L^1_{\rm loc})$ 理解。上述强紧性、标量结构与熵唯一性是关键，不能把同样的结论直接移植到任意守恒律系统。

% 第二次扩充（第6组）：非线性扩散、完全非线性及高阶模型。
% References: https://verso.mat.uam.es/~juanluis.vazquez/coursejlv.html
% https://www.tuwien.at/index.php?eID=dumpFile&f=1449534&t=f&token=1cf5a69a8dc1f8193e2a730445d8880b219abf4b
% https://bilman.github.io/teaching/pde/pde-lec-19/
\section{非线性扩散、Hessian方程与高阶模型}
\subsection{多孔介质方程与Barenblatt解}\index{多孔介质方程与Barenblatt解|hyperpage}
全空间非负扩散 $u_t=\Delta(u^m)$、$m>1$ 的质量守恒型自相似解为
\begin{gather*}
U(x,t)=t^{-\alpha}\left(C-k|x|^2t^{-2\beta}\right)_+^{1/(m-1)},\\
\alpha=\frac d{d(m-1)+2},\qquad
\beta=\frac1{d(m-1)+2},\qquad k=\frac{(m-1)\beta}{2m}.
\end{gather*}
其中 $C>0$ 由总质量确定，$U(t)\rightharpoonup M\delta_0$（$t\downarrow0$，测度意义）。支撑半径为 $\sqrt{C/k}\,t^\beta$。正值区的压力 $v=\frac m{m-1}u^{m-1}$ 满足
\[
v_t=(m-1)v\Delta v+|\nabla v|^2.
\]
退化性导致有限传播至真空；$0<m<1$ 为快扩散，不能沿用这一截断正部公式及其传播结论。
\subsection{Monge--Ampère方程与凸性分支}\index{Monge-Ampère方程与凸性分支|hyperpage}
对凸函数 $u$，典型方程为 $\det D^2u=f\ge0$。在严格凸光滑解附近，线性化为
\[
\left.\frac{\mathrm d}{\mathrm d\varepsilon}\det D^2(u+\varepsilon w)\right|_{0}
=\operatorname{cof}(D^2u):D^2w.
\]
Hessian 正定保证线性化椭圆；一致椭圆还要求 Hessian 特征值具有双侧正界。对 Borel 集 $E$，非光滑凸解用次梯度像定义 Alexandrov 测度 $\operatorname{MA}_u(E)=|\partial u(E)|$，光滑时 $\operatorname{MA}_u(E)=\int_E\det D^2u$。径向 $u(x)=h(r)$、$r>0$ 时，
\[
\det D^2u=h''(r)\left(\frac{h'(r)}r\right)^{d-1}.
\]
凸性是所选椭圆分支的条件；仅写出行列式等式不能保证椭圆性或边值问题唯一性。
\subsection{双调和方程与薄板边界}\index{双调和方程与薄板边界|hyperpage}
有界光滑域的夹持双调和问题为 $\Delta^2u=f$、$u=\partial_nu=0$。令 $H_0^2=\overline{C_c^\infty}^{\,H^2}$，$f\in H^{-2}$，其弱形式与能量为
\begin{gather*}
\int_\Omega\Delta u\Delta v\,\mathrm dx=\langle f,v\rangle
\quad(\forall v\in H_0^2),\\
J(u)=\frac12\int_\Omega(\Delta u)^2\,\mathrm dx-\langle f,u\rangle.
\end{gather*}
在 $H_0^2$ 上二阶能量强制，因此存在唯一弱解。边界 $u=\Delta u=0$ 则称 Navier 型，在光滑条件下可拆成两个 Dirichlet Poisson 问题：$-\Delta w=f$、$-\Delta u=w$。夹持条件与 Navier 条件不可互换；薄板自然边界还依赖所采用的弹性能量。
\subsection{KdV方程、孤立波与守恒量}\index{KdV方程、孤立波与守恒量|hyperpage}
采用标准化 $u_t+6uu_x+u_{xxx}=0$，对任意 $c>0$，有单孤立波
\[
u(x,t)=\frac c2\operatorname{sech}^2
\left(\frac{\sqrt c}{2}(x-ct-x_0)\right),\qquad x_0\in\mathbb R.
\]
其振幅为 $c/2$，宽度按 $c^{-1/2}$ 缩放。对足够光滑且衰减的解，下列积分守恒：
\[
\int u\,\mathrm dx,\qquad\int u^2\,\mathrm dx,\qquad
\int\left(\frac12u_x^2-u^3\right)\,\mathrm dx.
\]
线性 Airy 部分 $u_t+u_{xxx}=0$ 的角频率 Fourier 解为 $\widehat u(t,\xi)=e^{it\xi^3}\widehat u_0(\xi)$。非线性项系数或符号改变后，孤立波振幅及守恒能量中的常数须同步修改。

% 第二次扩充（第7组）：非局部算子。
% References: https://web.ma.utexas.edu/mediawiki/index.php/Fractional_Laplacian
% https://arxiv.org/abs/math/0608640 ; https://arxiv.org/abs/1808.05159
\section{分数阶算子、延拓与非局部边值}
\subsection{分数Laplacian的等价表示与能量}\index{分数Laplacian的等价表示与能量|hyperpage}
取角频率 Fourier 约定 $\widehat u(\xi)=\int e^{-ix\cdot\xi}u(x)\,\mathrm dx$，$0<s<1$。对 Schwartz 函数，
\begin{gather*}
\widehat{(-\Delta)^su}(\xi)=|\xi|^{2s}\widehat u(\xi),\\
(-\Delta)^su(x)=C_{d,s}\operatorname{p.v.}\int_{\mathbb R^d}
\frac{u(x)-u(y)}{|x-y|^{d+2s}}\,\mathrm dy,\qquad
C_{d,s}=\frac{4^s\Gamma(d/2+s)}{\pi^{d/2}|\Gamma(-s)|}.
\end{gather*}
对实值 $u$，对称化给出
\[
\int u(-\Delta)^su\,\mathrm dx
=\frac{C_{d,s}}2\iint\frac{|u(x)-u(y)|^2}{|x-y|^{d+2s}}\,\mathrm dx\,\mathrm dy
=\|(-\Delta)^{s/2}u\|_2^2.
\]
主值与归一化不可省略；若 Fourier 指数改成 $-2\pi ix\cdot\xi$，乘子随之变为 $(2\pi|\xi|)^{2s}$。
\subsection{热半群表示与Caffarelli--Silvestre延拓}\index{热半群表示与Caffarelli-Silvestre延拓|hyperpage}
对足够好的 $u$，有
\[
(-\Delta)^su=\frac1{\Gamma(-s)}\int_0^\infty
\bigl(e^{t\Delta}u-u\bigr)\frac{\mathrm dt}{t^{1+s}}.
\]
令 $U$ 为上半空间中有限加权能量的延拓，满足
\[
\operatorname{div}_{x,y}(y^{1-2s}\nabla U)=0,\qquad U(x,0)=u(x).
\]
其加权法导数与能量满足
\begin{gather*}
(-\Delta)^su=-\kappa_s\lim_{y\downarrow0}y^{1-2s}\partial_yU,\qquad
\kappa_s=\frac{2^{2s-1}\Gamma(s)}{\Gamma(1-s)},\\
\|(-\Delta)^{s/2}u\|_2^2
=\kappa_s\int_{\mathbb R^{d+1}_+}y^{1-2s}|\nabla U|^2\,\mathrm dx\,\mathrm dy.
\end{gather*}
对一般 $H^s$ 边值，法导数按弱意义解释；$s=1/2$ 时为通常调和延拓与 Dirichlet--Neumann 映射。
\subsection{分数热核、尺度与平滑化}\index{分数热核、尺度与平滑化|hyperpage}
全空间 $u_t+(-\Delta)^su=0$ 的核 $p_s$ 满足
\[
\widehat p_s(t,\xi)=e^{-t|\xi|^{2s}},\qquad
p_s(t,x)=t^{-d/(2s)}p_s(1,t^{-1/(2s)}x).
\]
对 $t>0$、$0<s<1$，它非负、总积分为一，并有双侧估计
\[
p_s(t,x)\asymp\frac{t}{(t^{1/(2s)}+|x|)^{d+2s}},
\]
比较常数仅依赖 $d,s$。因此对 $1\le p\le q\le\infty$，
\[
\|p_s(t)*f\|_q\le C t^{-\frac d{2s}(\frac1p-\frac1q)}\|f\|_p.
\]
此核具有代数尾，与 Gaussian 核的尾部不同；$s\uparrow1$ 时不能把固定 $s<1$ 的尾部比较常数直接当作一致常数使用。
\subsection{外部Dirichlet条件与非局部弱解}\index{外部Dirichlet条件与非局部弱解|hyperpage}
有界 Lipschitz 域 $\Omega$ 中，零外部条件要求 $u=0$ 于整个 $\mathbb R^d\setminus\Omega$。令
\[
V_s=\{u\in H^s(\mathbb R^d):u=0\ \text{a.e. 于 }\Omega^c\},
\]
\[
\mathcal E_s(u,v)=\frac{C_{d,s}}2\iint_{\mathbb R^d\times\mathbb R^d}
\frac{(u(x)-u(y))(v(x)-v(y))}{|x-y|^{d+2s}}\,\mathrm dx\,\mathrm dy.
\]
对 $f\in V_s^*$，解由 $\mathcal E_s(u,v)=\langle f,v\rangle$（所有 $v\in V_s$）唯一确定，亦极小化 $\frac12\mathcal E_s(u,u)-\langle f,u\rangle$。若 $f$ 对非负测试非负，则 $u\ge0$。这里只在域内施加方程，却必须给出整个外部的值；该限制型分数算子与 Dirichlet Laplace 特征展开定义的谱分数算子一般不同。

% 第二次扩充（第8组）：不可压流体。
% References: https://virtualmath1.stanford.edu/~ryzhik/notes-256B-24.pdf
% https://terrytao.wordpress.com/2018/10/02/254a-notes-2-weak-solutions-of-the-navier-stokes-equations/
\section{不可压流体的投影、弱解与延拓判据}
\subsection{Leray投影、压力与温和解}\index{Leray投影、压力与温和解|hyperpage}
在 $\mathbb R^d$ 上，Leray 投影到无散向量场的 Fourier 符号为
\[
\widehat{\mathbb P v}(\xi)=\left(I-\frac{\xi\otimes\xi}{|\xi|^2}\right)\widehat v(\xi)
\quad(\xi\ne0),\qquad\mathbb P=I-\nabla\Delta^{-1}\operatorname{div}.
\]
对无外力 Navier--Stokes，$\nu>0$、$\operatorname{div}u_0=0$，相应积分形式为
\[
u(t)=e^{\nu t\Delta}u_0-\int_0^te^{\nu(t-s)\Delta}
\mathbb P\operatorname{div}(u\otimes u)(s)\,\mathrm ds.
\]
压力满足 $-\Delta p=\sum_{i,j}\partial_i\partial_j(u_i u_j)$；在适当衰减与可积条件下 $p=\sum_{i,j}R_iR_j(u_i u_j)$，其中 $\widehat{R_if}=-i\xi_i|\xi|^{-1}\widehat f$。压力仍可加仅随时间的常数；有界域的投影和压力须结合边界条件重新定义。
\subsection{Biot--Savart公式与涡量恢复}\index{Biot-Savart公式与涡量恢复|hyperpage}
全空间适当衰减的无散速度可由涡量恢复。二维取 $z^\perp=(-z_2,z_1)$、$\omega=\partial_1u_2-\partial_2u_1$，有
\[
u(x)=\frac1{2\pi}\int_{\mathbb R^2}
\frac{(x-y)^\perp}{|x-y|^2}\omega(y)\,\mathrm dy.
\]
三维取 $\omega=\nabla\times u$，则
\[
u(x)=\frac1{4\pi}\int_{\mathbb R^3}
\frac{\omega(y)\times(x-y)}{|x-y|^3}\,\mathrm dy.
\]
导数按奇异积分理解，$\|\nabla u\|_p\le C_p\|\omega\|_p$（$1<p<\infty$）。在周期域还需指定平均速度，有边界时还需边界与可能的调和分量；不能仅凭涡量任意忽略这些自由度。
\subsection{Leray--Hopf能量不等式与弱强唯一性}\index{Leray-Hopf能量不等式与弱强唯一性|hyperpage}
无外力、全空间或三维周期域中，对无散 $u_0\in L^2$，Leray--Hopf 解属于
\[
u\in L^\infty(0,T;L^2)\cap L^2(0,T;H^1)
\]
（全空间中可等价记录 $L^2$ 范数与 $L^2$ 梯度）。其能量不等式为几乎每个 $t$，
\[
\frac12\|u(t)\|_2^2+\nu\int_0^t\|\nabla u(s)\|_2^2\,\mathrm ds
\le\frac12\|u_0\|_2^2.
\]
若 $v$ 为同方程的强解，且 $\int_0^T\|\nabla v\|_\infty<\infty$，相对能量估计给
\[
\|u(t)-v(t)\|_2^2\le\|u_0-v_0\|_2^2
\exp\left(2\int_0^t\|\nabla v(s)\|_\infty\,\mathrm ds\right).
\]
因此同初值弱解与该强解在其存在区间一致。能量不等式允许弱解的耗散，不能自动提升为每个弱解的能量等式或任意两个弱解的唯一性。
\subsection{Serrin条件、BKM判据与尺度}\index{Serrin条件、BKM判据与尺度|hyperpage}
三维无外力 Navier--Stokes 的尺度为 $u_\lambda(x,t)=\lambda u(\lambda x,\lambda^2t)$、$p_\lambda=\lambda^2p(\lambda x,\lambda^2t)$。相应速度范数满足
\[
\|u_\lambda\|_{L^r(0,T/\lambda^2;L^q)}
=\lambda^{1-2/r-3/q}\|u\|_{L^r(0,T;L^q)}.
\]
对已有的局部强解，Prodi--Serrin 型延拓条件为
\[
u\in L^r(0,T;L^q),\qquad\frac2r+\frac3q\le1,\qquad3<q\le\infty.
\]
该条件防止在 $T$ 发生强解奇性；$q=3,r=\infty$ 是需单独处理的端点。三维 Euler 在 $H^s$、$s>5/2$ 的局部强解有 Beale--Kato--Majda 判据：若最大正时间 $T_*<\infty$，必有
\[
\int_0^{T_*}\|\omega(t)\|_\infty\,\mathrm dt=\infty.
\]
同类涡量充分延拓条件也适用于光滑无外力 Navier--Stokes。条件性正则判据不能当作已由初值自动控制的全局先验界。

% References: https://virtualmath1.stanford.edu/~ryzhik/notes-256B-24.pdf
% https://web.ma.utexas.edu/mediawiki/index.php/Extremal_operator
% https://www.ocw.mit.edu/courses/18-354j-nonlinear-dynamics-ii-continuum-systems-spring-2015/6163837c5efa5d6ed969e535ec8f890c_MIT18_354JS15_lectureNotes.pdf
\section{非线性抛物模型与耗散}
\subsection{半线性热方程与Fujita指数}\index{半线性热方程与Fujita指数|hyperpage}
对全空间 $u_t=\Delta u+u^p$、$p>1$，非负解满足积分形式
\[
u(t)=e^{t\Delta}u_0+\int_0^te^{(t-s)\Delta}u(s)^p\,\mathrm ds.
\]
Fujita 临界指数为 $p_F=1+2/d$。对非零、非负、光滑紧支撑初值，$1<p\le p_F$ 时不存在全局有界经典解；$p>p_F$ 时足够小的 $L^1\cap L^\infty$ 初值有全局解，而大初值仍可爆破。此处源项为正幂增长，吸收方程 $u_t=\Delta u-u^p$ 不遵循同一爆破结论。
\subsection{多孔介质的Lp与压力耗散}\index{多孔介质的Lp与压力耗散|hyperpage}
对 $u_t=\Delta(u^m)$、$m>1$ 的足够好的非负全空间解，$r>1$ 时
\[
\frac1r\frac{\mathrm d}{\mathrm dt}\int u^r
=-\frac{4m(r-1)}{(m+r-1)^2}\int\left|\nabla u^{(m+r-1)/2}\right|^2.
\]
压力 $v=mu^{m-1}/(m-1)$ 给出 $\frac{\mathrm d}{\mathrm dt}\frac1{m-1}\int u^m=-\int u|\nabla v|^2$。对弱解通常以相应耗散不等式及逼近理解，须控制正值区、可积性和无穷远通量。
\subsection{Allen--Cahn能量下降}\index{Allen-Cahn能量下降|hyperpage}
令 $W(s)=(1-s^2)^2/4$，取 $u_t=\Delta u-\varepsilon^{-2}W'(u)$。周期域或齐次 Neumann 边界下，光滑解满足
\[
E_\varepsilon(u)=\int\left(\frac\varepsilon2|\nabla u|^2+\frac{W(u)}\varepsilon\right),\qquad
\frac{\mathrm d}{\mathrm dt}E_\varepsilon(u)=-\varepsilon\int|u_t|^2.
\]
该流一般不守恒总质量。若改变时间尺度，右端系数也会改变；相场能量的周长极限见\hyperref[chap:calculus-of-variations]{变分法}。
\subsection{Cahn--Hilliard的质量与能量}\index{Cahn-Hilliard的质量与能量|hyperpage}
取 $u_t=\Delta\mu$、$\mu=-\varepsilon\Delta u+\varepsilon^{-1}W'(u)$，在周期域或 $\partial_nu=\partial_n\mu=0$ 下，
\[
\frac{\mathrm d}{\mathrm dt}\int u=0,\qquad
\frac{\mathrm d}{\mathrm dt}E_\varepsilon(u)=-\int|\nabla\mu|^2.
\]
这是同一相场能量的负 Sobolev 度量梯度流；与 Allen--Cahn 相比，度量变化使演化成为四阶并保持质量。弱解的能量表述需相应函数空间。
\section{完全非线性算子与椭圆分支}
\subsection{Pucci极值算子}\index{Pucci极值算子|hyperpage}
对称矩阵 $X$ 的特征值为 $e_i$，$0<\lambda\le\Lambda$，定义
\begin{align*}
\mathcal M^+(X)&=\Lambda\sum_{e_i>0}e_i+\lambda\sum_{e_i<0}e_i,\\
\mathcal M^-(X)&=\lambda\sum_{e_i>0}e_i+\Lambda\sum_{e_i<0}e_i.
\end{align*}
它们分别是 $\lambda I\preceq A\preceq\Lambda I$ 下 $\operatorname{tr}(AX)$ 的上、下确界。采用关于 Hessian 递增的符号约定，一致椭圆 $F$ 满足 $\mathcal M^-(X-Y)\le F(X)-F(Y)\le\mathcal M^+(X-Y)$。
\subsection{Legendre对偶下的Hessian方程}\index{Legendre对偶下的Hessian方程|hyperpage}
若 $u$ 光滑严格凸、梯度为微分同胚，则对 $y=\nabla u(x)$、$u^*$ 为 Legendre 共轭，
\[
\nabla u^*(y)=x,\qquad D^2u^*(y)=(D^2u(x))^{-1}.
\]
故 $\det D^2u=f(x)>0$ 变为 $\det D^2u^*(y)=1/f(\nabla u^*(y))$。对偶定义域是原梯度像；不应忽略定义域而把该变换当作同域上的简单取倒数。
\subsection{k-Hessian与Newton张量}\index{k-Hessian与Newton张量|hyperpage}
设 $1\le k\le d$。记 $S_k(D^2u)$ 为 Hessian 特征值的第 $k$ 个基本对称多项式，$T_{k-1}^{ij}=\partial S_k/\partial u_{ij}$。对光滑 $u$，
\[
\partial_iT_{k-1}^{ij}=0,\qquad
S_k(D^2u)=\frac1k\operatorname{div}(T_{k-1}(D^2u)\nabla u).
\]
$k$-可容许分支由特征值位于 $\Gamma_k=\{S_j>0:1\le j\le k\}$ 给出。在该分支内线性化椭圆；离开分支仅有多项式等式，不能继续援用同样的椭圆理论。
\subsection{无穷Laplacian与p趋于无穷}\index{无穷Laplacian与p趋于无穷|hyperpage}
定义 $\Delta_\infty u=\sum_{i,j}u_i u_j u_{ij}$。当 $\nabla u\ne0$ 时，
\[
\Delta_pu=|\nabla u|^{p-2}
\left(\Delta u+(p-2)\frac{\Delta_\infty u}{|\nabla u|^2}\right),\qquad
\Delta_\infty u=\tfrac12\nabla u\cdot\nabla|\nabla u|^2.
\]
归一化后的极限关联绝对极小 Lipschitz 延拓及 $\Delta_\infty u=0$ 的黏性解。梯度为零处不能直接除以 $|\nabla u|^2$，需按黏性测试的相应定义处理。
\section{流体基本解、输运与二维结构}
\subsection{三维Stokes基本解}\index{三维Stokes基本解|hyperpage}
对 $-\nu\Delta u+\nabla p=F\delta_0$、$\operatorname{div}u=0$，$\nu>0$，全空间基本解为
\[
u_i(x)=\sum_jG_{ij}(x)F_j,\qquad
G_{ij}(x)=\frac1{8\pi\nu}\left(\frac{\delta_{ij}}{|x|}+\frac{x_ix_j}{|x|^3}\right),
\]
\[
p(x)=\frac{F\cdot x}{4\pi|x|^3}.
\]
此为 Stokeslet/Oseen 张量，源项等式按分布成立；二维存在对数行为及 Stokes 悖论，不能把 $d=3$ 核直接改写下标使用。
\subsection{Kelvin环量与Euler螺度}\index{Kelvin环量与Euler螺度|hyperpage}
对光滑无外力不可压 Euler 解，随流体移动的闭曲线 $C(t)$ 满足
\[
\frac{\mathrm d}{\mathrm dt}\oint_{C(t)}u\cdot\mathrm dx=0.
\]
三维充分衰减或周期解还守恒螺度 $\mathcal H=\int u\cdot\omega\,\mathrm dx$，$\omega=\nabla\times u$。黏性、非保守外力或不适当边界通量会改变这些结论，不能与 Navier--Stokes 无条件混用。
\subsection{Yudovich涡量与对数Lipschitz控制}\index{Yudovich涡量与对数Lipschitz控制|hyperpage}
二维全空间中，$\omega\in L^1\cap L^\infty$ 的 Biot--Savart 速度满足
\begin{gather*}
\|u\|_\infty\le C\|\omega\|_1^{1/2}\|\omega\|_\infty^{1/2},\\
|u(x)-u(y)|\le C(\|\omega\|_1+\|\omega\|_\infty)
|x-y|(1+|\log|x-y||),\quad0<|x-y|\le1.
\end{gather*}
Euler 输运保持相应涡量范数；对数 Lipschitz 模满足 Osgood 唯一性积分条件。配合速度在无穷远的归一化，可构造该有界涡量类中的全局唯一解；这不是对所有粗糙 Euler 弱解的唯一性断言。
\subsection{连续性方程的Jacobian公式}\index{连续性方程的Jacobian公式|hyperpage}
对足够光滑且保证全局流存在的速度 $b$，$X_t(t,a)=b(t,X(t,a))$、$X(0,a)=a$，有
\[
J(t,a)=\det D_aX(t,a)
=\exp\left(\int_0^t\operatorname{div}b(s,X(s,a))\,\mathrm ds\right).
\]
连续性方程 $\rho_t+\operatorname{div}(\rho b)=0$ 给出 $\rho(t,X(t,a))J(t,a)=\rho_0(a)$；而纯输运方程给出沿流函数值不变。两式只有在 $\operatorname{div}b=0$ 等情形才具有相同的密度表达。
\section{边界谱、惩罚与移动界面}
\subsection{Dirichlet--Neumann映射与Steklov谱}\index{Dirichlet-Neumann映射与Steklov谱|hyperpage}
光滑有界连通域上，$Hg$ 为边值 $g\in H^{1/2}(\partial\Omega)$ 的调和延拓，定义 $\Lambda g=\partial_nHg\in H^{-1/2}$。有
\[
\langle\Lambda g,h\rangle=\int_\Omega\nabla Hg\cdot\nabla Hh,\qquad
\langle\Lambda g,g\rangle\ge0.
\]
Steklov 问题 $\Delta u=0$、$\partial_nu=\sigma u$ 是该边界算子的谱问题；常数给零特征值，其余特征值非负并离散趋于无穷。与 Dirichlet Laplace 谱相比，归一化质量位于边界。
\subsection{Robin惩罚趋于Dirichlet}\index{Robin惩罚趋于Dirichlet|hyperpage}
对 $-\Delta u_\beta+u_\beta=f$、$\partial_nu_\beta+\beta u_\beta=0$，$\beta>0$，相应能量为
\[
J_\beta(u)=\frac12\int_\Omega(|\nabla u|^2+u^2)
+\frac\beta2\int_{\partial\Omega}|Tu|^2-\int_\Omega fu.
\]
有界 Lipschitz 域及 $f\in L^2$ 下，$\|Tu_\beta\|_{L^2(\partial\Omega)}\le C\beta^{-1/2}$，且 $u_\beta\to u_D$ 于 $H^1$，其中 $u_D$ 为零 Dirichlet 解。边界惩罚的正号与增加的零阶项保证这里的统一强制性。
\subsection{波动的吸收边界与能量通量}\index{波动的吸收边界与能量通量|hyperpage}
对齐次波动方程，在光滑边界施加 $u_t+c\partial_nu=0$，则
\[
\frac{\mathrm d}{\mathrm dt}\frac12\int_\Omega(u_t^2+c^2|\nabla u|^2)
=-c\int_{\partial\Omega}u_t^2\le0.
\]
一维右端条件为 $u_t+cu_x=0$，左端为 $u_t-cu_x=0$，可精确吸收对应外行波。多维该局部阻抗条件一般只是吸收近似，不应声称对任意入射角完全无反射。
\subsection{Stefan界面的守恒跳跃}\index{Stefan界面的守恒跳跃|hyperpage}
若焓 $H$ 与热流 $q$ 满足 $H_t+\operatorname{div}q=0$，跨光滑移动界面，取法向从负侧指向正侧、$[H]=H^+-H^-$，则
\[
V_n[H]=[q]\cdot n.
\]
在两相 Fourier 导热 $q^\pm=-\kappa_\pm\nabla T^\pm$ 下，$V_n[H]=[-\kappa\partial_nT]$，并通常附加界面温度等于相变温度。潜热跳跃的符号由两侧约定决定；先固定 $n$ 与跳跃定义再使用速度公式。

\section{电磁、弹性与动理学方程补充}
下列守恒恒等式按光滑解陈述，边界或无穷远通量需满足相应条件。弹性基础见\cite{variation}，气体动理学见\cite{kinetic}。
\subsection{Maxwell系统与Poynting恒等式}\index{Maxwell系统与Poynting恒等式|hyperpage}
常数介电率 $\varepsilon>0$、磁导率 $\mu>0$ 下，三维 Maxwell 演化方程为
\[
\varepsilon E_t=\nabla\times H-J,\qquad
\mu H_t=-\nabla\times E,
\]
并有 $\operatorname{div}(\varepsilon E)=\rho$、$\operatorname{div}(\mu H)=0$。能量密度满足
\[
\partial_t\frac{\varepsilon|E|^2+\mu|H|^2}{2}
+\operatorname{div}(E\times H)=-J\cdot E.
\]
在周期域或完全导体边界 $n\times E=0$ 下，边界能流消失；若 $J=\sigma E$、$\sigma\ge0$，总场能非增。电荷连续性保证散度约束被演化保持。
\subsection{Navier--Lamé弹性系统与Korn控制}\index{Navier-Lamé弹性系统与Korn控制|hyperpage}
记 $e(u)=(\nabla u+\nabla u^{\mathsf T})/2$，常系数各向同性应力为 $\sigma(u)=2\mu e(u)+\lambda\operatorname{div}u\,I$。平衡式为
\[
-\mu\Delta u-(\lambda+\mu)\nabla\operatorname{div}u=f.
\]
取 $\mu>0$、$\lambda+2\mu/d>0$，弹性能量正定。齐次固定边界时 Korn 不等式 $\|\nabla u\|_2\le C\|e(u)\|_2$ 排除刚体运动；自由边界时需商去相应有限维核。
\subsection{Boltzmann碰撞不变量与H定理}\index{Boltzmann碰撞不变量与H定理|hyperpage}
弹性碰撞算子写作 $Q(f,f)(v)=\iint B(f'f_*'-ff_*)\,\mathrm d\sigma\mathrm dv_*$，其中
\[
v'=\frac{v+v_*}{2}+\frac{|v-v_*|}{2}\sigma,\qquad
v_*'=\frac{v+v_*}{2}-\frac{|v-v_*|}{2}\sigma.
\]
对具有通常微观可逆对称性的非负碰撞核、足够可积的正密度，
\[
\int Q(f,f)\begin{pmatrix}1\\v\\|v|^2\end{pmatrix}\,\mathrm dv=0,
\qquad \int Q(f,f)\log f\,\mathrm dv\le0.
\]
非退化碰撞下平衡态为 Maxwell 密度 $\exp(a+b\cdot v-c|v|^2)$、$c>0$。弱解表述需保证各积分或相应耗散不等式有意义。
\newpage
\subsection{Vlasov--Poisson质量与场能}\index{Vlasov-Poisson质量与场能|hyperpage}
单位体积周期空间中，取排斥号约定
\[
f_t+v\cdot\nabla_xf+E\cdot\nabla_vf=0,\qquad
E=-\nabla\phi,\quad-\Delta\phi=\rho-\bar\rho,\quad
\rho=\int f\,\mathrm dv.
\]
取零均值势，$\bar\rho=\int\rho\,\mathrm dx$。光滑解在速度无穷远充分衰减时，质量与能量
\[
\mathcal E=\frac12\iint |v|^2f\,\mathrm dv\mathrm dx
+\frac12\int|E|^2\,\mathrm dx
\]
守恒；吸引型系统须改变 Poisson 与场能符号。
\clearpage
\end{document}
